\documentclass[11pt]{article}

\usepackage[margin=1in]{geometry}
\usepackage{amsmath,amssymb,bm,mathtools}
\usepackage{booktabs}
\usepackage{longtable,array}
\usepackage{enumitem}
\usepackage{microtype}
\usepackage{xcolor}
\usepackage{hyperref}

\hypersetup{colorlinks=true,linkcolor=blue!50!black,urlcolor=blue!60!black}

\newcommand{\G}{\mathcal{G}}
\newcommand{\Gt}{\mathcal{G}_{\theta}}
\newcommand{\R}{\mathcal{R}_{\theta}}
\newcommand{\dd}{\mathop{}\!\mathrm{d}}
\newcommand{\norm}[1]{\left\lVert #1\right\rVert}
\newcommand{\inner}[2]{\left\langle #1,#2\right\rangle}
\newcommand{\Deta}{D_{\eta}}
\allowdisplaybreaks

\title{Diagnosing and Correcting NaN Cascades in Learned\\
Dirichlet--Neumann Operator Rollouts}
\author{Internal technical report}
\date{9 July 2026}

\begin{document}
\maketitle
\tableofcontents
\newpage

\begin{abstract}
Accurate one-step approximation of a water-wave Dirichlet--Neumann operator
(DNO) did not prevent rare catastrophic failures in long symplectic rollouts.
Frame-aligned trajectory analysis shows why: the learned operator is accurate
on exact solution states but has an incorrect derivative transverse to the
solution manifold.  It consequently violates the Hadamard shape-derivative
identity, so the two equations advanced by the Gauss--Legendre integrator are
not the canonical derivatives of one learned Hamiltonian.  Small state errors
then excite spurious same-sign spectral sidebands, quadratic aliasing amplifies
the already established cascade, and the implicit stage iteration eventually
loses contractivity.  This report summarizes the experimental path that ruled
out data corruption, time-step refinement, spectral caps, hard cutoffs, and
targeted fine-tuning as primary cures.  It motivates a clean retraining of
\(G_0+G_1+\R\), where the analytic first-order Craig--Sulem term is exact,
\(\R=O(\eta^2)\) is self-adjoint, and a randomized Hadamard loss directly
constrains the missing shape derivative.  Particular care is taken to separate
what has been directly measured, what follows mathematically from the measured
facts, and what remains a hypothesis to be tested by the new model.
\end{abstract}

\section{Reader's guide and executive conclusions}

The project accumulated many experiments because the terminal symptom---a NaN
inside an implicit time step---admits several superficially plausible
explanations.  A reader should not be expected to remember the experiment
labels.  The labels C1--C16 are therefore treated here only as reproducibility
handles.  The scientific argument is organized around five questions:
\begin{enumerate}[label=\arabic*.]
  \item Is the reference data corrupt or missing the failing regime?
  \item Does the learned DNO first make an error on an exact physical state, or
        does it react incorrectly to the small off-manifold state error created
        by its own rollout?
  \item Is the cascade initiated by the DNO, by pseudospectral aliasing, or by
        nonconvergence of the implicit Gauss--Legendre stages?
  \item Which exact property of a physical DNO is absent from the learned map?
  \item What is the smallest architectural and loss-level correction that
        restores that property without suppressing legitimate high-frequency
        dynamics?
\end{enumerate}

The present answers, with their evidentiary status, are:
\begin{description}[style=nextline,leftmargin=0pt]
\item[Directly measured:]
The failing models can have relative DNO value error of only
\(2\times10^{-4}\) on the exact state while their response to the difference
between the exact and surrogate states grows to a secant gain of 7--8.  The
induced response first grows in \(k=32,\ldots,64\), then contaminates
\(k=64,\ldots,128\), and only afterward makes the quadratic right-hand side
strongly aliased and the fixed-point stage solve noncontractive.

\item[Directly measured structural defect:]
At the same exact states, the learned operator violates the DNO Hadamard
shape-derivative identity by relative \(O(1)\), whereas the order-six reference
implementation satisfies the identical discrete test to approximately
\(10^{-5}\) or better.  This discrepancy is far larger than the DNO value
error and is already present before visible blow-up.

\item[Mathematical consequence:]
If a learned DNO is self-adjoint in \(\xi\) but fails the shape identity, then
the implemented formulas for \(\eta_t\) and \(\xi_t\) are not the two
variational derivatives of one learned Hamiltonian.  A symplectic integrator
cannot restore a Hamiltonian structure absent from the vector field it is
given.

\item[Strong mechanistic evidence:]
The first Craig--Sulem term has a same-sign Fourier null form that cancels the
specific high-mode sidebands observed in the failures.  The learned operator
does not preserve this cancellation.  Exact \(G_0+G_1\) gives 0/32 Tanaka
NaNs but is not accurate enough at final time, strongly motivating it as a
stable backbone rather than a complete surrogate.

\item[Hypothesis under test:]
A full-capacity model of the form
\(G_0+G_1+O(\eta^2)\), trained with a direct Hadamard-consistency loss, should
retain the stable null form while learning the small higher-order correction.
This is C16.  Its structural invariants and short training smoke have passed;
the decisive long-rollout evaluation has not yet been performed.
\end{description}

This hierarchy matters.  We have compelling evidence for the diagnosis, but
we do not yet have evidence that C16 reaches the operational target of zero
NaNs.  A negative C16 result would not erase the measured shape defect; it
would mean that the particular parameterization, loss estimator, training
budget, or some additional missing identity remains inadequate.

\section{Problem and numerical setting}

We consider a periodic one-dimensional horizontal domain of length \(2\pi\),
with free surface \(y=\eta(x,t)\), constant depth \(h\), and an irrotational
velocity potential \(\phi\) in the fluid.  Its trace at the free surface is
\(\xi(x,t)=\phi(x,\eta(x,t),t)\).  The Dirichlet--Neumann operator is the map
from this trace to the non-unit normal derivative,
\begin{equation}
 G(\eta;h)\xi
 =\left.\bigl(\phi_y-\eta_x\phi_x\bigr)\right|_{y=\eta}.
 \label{eq:dno-definition}
\end{equation}
The definition incorporates the solution of Laplace's equation in the entire
fluid domain.  It is therefore the expensive nonlocal component that the
surrogate is intended to replace.

For surface elevation \(\eta\) and surface potential \(\xi\), Zakharov's
canonical water-wave equations without surface tension may be written
\begin{align}
  \eta_t &= G(\eta)\xi, \\
  \xi_t &= -g\eta
  -\frac{1}{2}\xi_x^2
  +\frac{1}{2}\frac{\left(G(\eta)\xi+\eta_x\xi_x\right)^2}
                         {1+\eta_x^2}.
 \label{eq:zakharov-geometric}
\end{align}
They are Hamilton's equations for
\begin{equation}
 H(\eta,\xi)
 =\frac12\int_0^{2\pi}\xi G(\eta)\xi\,\dd x
  +\frac{g}{2}\int_0^{2\pi}\eta^2\,\dd x,
 \qquad
 \eta_t=\frac{\delta H}{\delta\xi},\quad
 \xi_t=-\frac{\delta H}{\delta\eta}.
 \label{eq:hamiltonian}
\end{equation}
This equivalence uses both self-adjointness of \(G(\eta)\) in its potential
argument and its precise derivative with respect to the surface.  Preserving
only the former does not suffice.

The production surrogate replaces the expensive DNO \(G\) by a learned
operator \(\Gt\), while retaining the geometric formula for \(\xi_t\).  Time
integration uses a two-stage Gauss--Legendre method with an integrating factor
and four fixed-point iterations per step.  The DNO model is linear in \(\xi\)
and built from Craig--Sulem-shaped multiplier sandwiches.

\subsection{Reference and surrogate operators}

The reference DNO is a truncated Craig--Sulem series evaluated through order
six with substantial Fourier zero padding.  It is not an exact continuum
oracle, but its convergence, self-adjointness, shape-identity residual, and
long-rollout conservation are all sufficiently accurate to distinguish the
effects discussed below.  In particular, the reference shape defect is
\(10^{-5}\) or smaller at the failing states, versus \(O(1)\) for the learned
operator.

The principal baseline is an eight-block Craig--Sulem DNO network.  Each block
is linear in \(\xi\): a depth-conditioned Fourier multiplier is applied to
\(\xi\), multiplied pointwise by learned features of \(\eta\), and passed
through another Fourier multiplier.  This is a much stronger inductive bias
than a generic FNO and produced excellent one-step validation error.  Later
checkpoints tied the input and output multipliers so that the learned residual
is self-adjoint in \(\xi\).  Nothing in the original architecture, however,
forces the correct derivative with respect to \(\eta\).

\subsection{The integrating-factor GL2 scheme}

Following Xu and Guyenne's spectral Hamiltonian method~\cite{xu2009}, split the
semi-discrete vector field as
\begin{equation}
 u_t=A u+N(u),\qquad u=(\eta,\xi)^\mathsf{T},qquad
 A=\begin{pmatrix}0&G_0\\-g&0\end{pmatrix},
 \label{eq:split}
\end{equation}
where \(G_0=|D|\tanh(h|D|)\).  With \(v(t)=e^{-tA}u(t)\), the stiff linear
dispersion is treated exactly and the transformed equation is advanced with
the two-stage, fourth-order Gauss--Legendre Runge--Kutta method.  Its Butcher
coefficients are
\begin{equation}
 a_{11}=a_{22}=\frac14,\quad
 a_{12}=\frac14-\frac{\sqrt3}{6},\quad
 a_{21}=\frac14+\frac{\sqrt3}{6},\quad
 b_1=b_2=\frac12,
 \label{eq:gl2-coefficients}
\end{equation}
up to interchange of the two stage labels.  The coupled stage equations are
solved here with four fixed-point iterations.

Three logically distinct statements must not be conflated:
\begin{enumerate}
  \item GL2 is a symplectic Runge--Kutta method for a canonical Hamiltonian
        vector field.
  \item A finite number of Picard iterations only approximately solves the
        implicit stage equations; loss of contractivity can expose a bad state.
  \item If the supplied surrogate vector field is not Hamiltonian, exact
        solution of the GL2 stages still cannot confer the missing Hamiltonian
        geometry on that field.
\end{enumerate}
The investigation below identifies item 3 as the initiating issue and item 2
as a terminal failure mechanism.

\subsection{What counts as success}

The practical target is stricter than small one-step validation loss: across
the 32-case Tanaka test and the 31 truth-valid Benjamin--Feir cases, the rollout
must produce no nonfinite trajectories while retaining competitive median and
tail error at final time.

We distinguish four metrics:
\begin{description}[style=nextline,leftmargin=0pt]
\item[One-step DNO error]
The relative error in \(G(\eta)\xi\) when the model is evaluated on an exact
reference state.  This measures interpolation on the sampled truth manifold.
\item[NaN rate]
The fraction of truth-valid initial conditions for which any surrogate state
becomes nonfinite.  This is the primary hard stability endpoint.
\item[Divergence rate]
The fraction whose final relative state error exceeds the evaluation threshold,
including but not limited to nonfinite cases.  A finite but physically useless
trajectory is not counted as a success.
\item[Final-time error distribution]
The median and 95th percentile relative \(L^2\) error in \(\eta\).  Reporting
only the median can hide a cap or filter that protects easy cases while
producing a catastrophic tail.
\end{description}

The rollout harness stores the integrator state and spectral operations in
double precision while leaving the neural network in its trained precision.
This matters: an earlier single-precision harness imposed an artificial
long-time error floor.  Truth-invalid initial conditions are excluded using
the reference solver's own nonfiniteness or excessive Hamiltonian drift, so a
surrogate is not penalized for a failed reference trajectory.

\section{What the experimental sequence established}

Table~\ref{tab:history} condenses the main branches recorded in
\texttt{EXPERIMENTS.md}.  The sequence is informative because several
apparently plausible interventions improved an intermediate diagnostic without
eliminating the failure population.

\subsection{Baseline numbers and checkpoint conventions}

The principal comparison points are listed in
Table~\ref{tab:checkpoint-baselines}.  ``v10'' denotes a fresh full-capacity v9
baseline.  C1 is a light stage-tangent fine-tune.  The first attempted C2 was
invalid because a tied-multiplier checkpoint was silently loaded into an
untied template; an architecture guard was added, and only the corrected C2
should be interpreted scientifically.  C10 is the strongest stage-gain variant
and the best pre-C16 NaN count.  These distinctions explain why an isolated
statement such as ``C2 failed'' can be misleading without a run path and
checkpoint choice.

\begin{table}[ht]
\centering
\small
\begin{tabular}{lrrrrrr}
\toprule
Model & Tanaka NaN & Tanaka div. & Tanaka med. & BF NaN & BF div. & BF med. \\
\midrule
v10 baseline & 4/32 & 5/32 & 0.0123 & 2/31 & 2/31 & 0.00902 \\
C1 stage match & 3/32 & 3/32 & 0.0151 & 0/31 & 0/31 & 0.00661 \\
Corrected C2 final & 3/32 & 4/32 & 0.00907 & 0/31 & 0/31 & 0.00735 \\
C10 final & 2/32 & 3/32 & 0.0138 & 0/31 & 0/31 & 0.00971 \\
\bottomrule
\end{tabular}
\caption{Primary pre-C16 comparisons under the double-precision harness.  The
Benjamin--Feir denominator excludes one truth-invalid case.}
\label{tab:checkpoint-baselines}
\end{table}

For a fine-tune driven by an auxiliary rollout or tangent loss, ``best
validation'' and ``best physical checkpoint'' need not coincide.  The ordinary
validation score measures the pointwise supervised objective and can select an
early epoch before the auxiliary constraint has acted.  This was demonstrated
by the corrected C2 run: the final checkpoint had the meaningful regularized
behavior, whereas the initially evaluated best-validation checkpoint from the
misconfigured run was unusable.  For C16, final/latest is therefore the primary
rollout checkpoint and best-validation is a secondary control.

\subsection{Phase I: establish a trustworthy evaluation}

Before asking why a model fails, we removed two ways the evaluation itself
could lie.
\begin{enumerate}
  \item A double-precision rollout harness reduced a spurious long-time error
        floor by factors of roughly 2.5 on Tanaka and 20 or more on several
        Benjamin--Feir regimes.  An order-six oracle displayed the same
        single-precision floor, proving it was not neural-model error.
  \item Reference trajectories with nonfinite states or Hamiltonian drift
        above the validity threshold were gated out.  For example, one
        Benjamin--Feir initial condition failed in the truth solver while the
        surrogate remained finite.  Counting it against the surrogate would
        invert the conclusion.
\end{enumerate}
The residual Tanaka NaNs survive both corrections and are therefore genuine
surrogate failures.

\subsection{Phase II: test the data-coverage explanation}

Additional steep-Tanaka data in v9 improved regimes that genuinely had a
coverage gap, including the Tanaka \(g_1\) family, but left 4/32 deep-water
Tanaka \(g_0\) NaNs.  Random and source-stratified scans found no nonfinite base
rows.  Earlier sign-flip heuristics were not treated as proofs of corruption;
the more decisive test came later: the learned DNO remained highly accurate
when evaluated directly on the exact trajectories that would eventually fail
under rollout.

Two forms of hard-negative training were then explored.  Model-harvested
pre-cascade states are circular: the inputs already contain the old model's
spectral contamination, and their approximate targets can inherit it.  These
packs worsened performance.  Clean truth-generated perturbations avoid that
specific flaw and are valid as local learnability probes, but they constrain a
finite list of neighborhoods rather than the universal derivative identity.
The clean C13 pack trained successfully, but its evaluation was intentionally
deprioritized once the structural diagnosis became available.  Thus the
conclusion is not ``additional data can never help''; it is that ordinary
pointwise labels are an inefficient and non-general remedy for an operator
derivative defect.

\subsection{Phase III: test the integrator explanation}

The GL2 stage residual becomes large immediately before nonfiniteness, making
time-step failure an attractive primary hypothesis.  Substep refinements and
halving trials, however, moved the failure time only modestly.  A matched
2-by-2 comparison---baseline versus C1, with and without Picard-residual
containment---left the NaN counts unchanged in every arm.  The containment
logic was useful operationally because it isolated bad members of a batch, but
it did not rescue the affected trajectories.

This phase did reveal a real training signal.  C1's stage-tangent loss cured
one of four baseline Tanaka NaNs and both Benjamin--Feir NaNs.  Later C10,
targeting the production-filtered \(k=64,\ldots,128\) stage gain, reduced
Tanaka NaNs to 2/32.  Hence sensitivity control is relevant.  The failure to
reach zero suggested that the indirect stage proxy was not enforcing the
defining differential relation strongly or accurately enough.

\subsection{Phase IV: suppress the observed spectrum}

Hard learned-block cutoffs, soft high-band penalties, run-time high-band
limiters, and a checkpoint-carrying output cap all attacked the visible
spectral symptom.  None reached zero NaNs with acceptable tails.  The strongest
soft penalty left its own high-band diagnostic essentially unchanged.  The
C10 run-time limiter retained both Tanaka NaNs and worsened the Benjamin--Feir
p95.  The C14 hard output projection retained 3/32 Tanaka NaNs, produced six
divergent Tanaka cases, and raised Tanaka p95 to 1.25.

The cap's algebra explains its failure.  It imposed
\begin{equation}
 \norm{g_{\mathrm{hi}}}
 \leq \max\!\left(5,0.1\norm{g_{\mathrm{lo}}}\right).
 \label{eq:old-cap}
\end{equation}
The dangerous \(k=16,\ldots,32\) content belonged to the unconstrained
``low'' group, growth of that group increased the permitted high-band norm,
and the absolute floor was inactive on some failing states.  This is a moving
relative envelope, not an absolute stability bound.

Finally, ablations showed why editing a trained checkpoint is hazardous.
Block 7 appeared individually amplifying, yet zeroing it changed the Tanaka NaN
count from 4/32 to 28/32 and the median error by a factor of 62.  Its
anti-\(|k|\) response was compensating the other blocks.  Likewise, removing a
large derivative feature from an already trained model breaks learned
cancellations and does not tell us whether a clean architecture should expose
that feature differently.

\begin{table}[ht]
\centering
\small
\begin{tabular}{p{0.18\linewidth}p{0.50\linewidth}p{0.23\linewidth}}
\toprule
Experiment family & Main evidence & Conclusion \\
\midrule
Clean v9 retraining & Excellent one-step error, yet the baseline retained
4/32 Tanaka NaNs.  Additional steep-Tanaka data cured a separate \(g_1\)
coverage failure but not the deep-water \(g_0\) failures. & Pointwise coverage
is not the principal remaining defect. \\

GL2 stage regularization & The best corrected fine-tune reduced Tanaka NaNs
from 4/32 to 3/32 and Benjamin--Feir NaNs to zero, but did not reach the
zero-NaN target and traded easy-case accuracy. & Directionally useful, but an
indirect and incompletely matched constraint. \\

Step residual checks and halving & Picard residual checks detected terminal
noncontractivity.  Step refinement and attempted halving moved failure times
only slightly; per-sample containment did not change the NaN count. & The
integrator exposes the failure after the model state is already contaminated. \\

Hard block cutoffs and output caps & Hard \(k\)-cuts and high-band envelopes
reduced selected spectral symptoms but produced poor error tails or retained
NaNs.  The one-block capped model reached 3/32 Tanaka NaNs and severe p95
regression. & Static culling removes physical response and leaves loopholes in
the growing mid-band. \\

Targeted hard negatives & Model-harvested pre-cascade fine-tunes were harmful;
clean truth-targeted packs did not provide a general or journal-explainable
mechanism. & Do not substitute case selection for an operator identity. \\

Exact \(G_0+G_1\) & A full 32-case CPU rollout gave 0/32 NaNs and 0/32
divergences, but final \(\eta\) error remained too large (median 0.207, p95
0.410). & The analytic null structure is stable, but needs an accurate
higher-order residual. \\
\bottomrule
\end{tabular}
\caption{Distilled experimental chronology.}
\label{tab:history}
\end{table}

Two broader lessons emerged.  First, deleting an apparently dangerous learned
block made the model dramatically worse (28/32 rather than 4/32 Tanaka NaNs),
showing that the trained eight-block representation relies on fragile
inter-block cancellation.  Second, increasing or decreasing parameter count
did not track rollout stability.  The missing ingredient is a structural
identity, not capacity alone.

\section{Frame-by-frame failure anatomy}

Let \(z_n^*=(\eta_n^*,\xi_n^*)\) be an exact truth state and \(z_n\) the
corresponding surrogate-rollout state.  Separate the model error into the
truth-manifold bias and the response induced by rollout displacement:
\begin{align}
 b_n &= \Gt(z_n^*)-G(z_n^*), \\
 i_n &= \Gt(z_n)-\Gt(z_n^*), \\
 q_n &= \frac{\norm{i_n}_2}{\norm{z_n-z_n^*}_2}.
\end{align}
The eventual failures remain accurate on truth states: median one-step
relative \(G\xi\) errors are \(1.87\times10^{-4}\) and
\(1.72\times10^{-4}\) for representative failing cases 5 and 11.  The
transverse response, however, grows by four orders of magnitude relative to
that bias.

\begin{table}[ht]
\centering
\small
\begin{tabular}{rrrrrr}
\toprule
Case & Time & State RMS error & Bias RMS & Induced RMS & Secant gain \\
\midrule
5  & 0.8   & $1.62\times10^{-6}$ & $5.12\times10^{-6}$ & $5.56\times10^{-6}$ & 3.43 \\
5  & 58.4  & $2.08\times10^{-3}$ & $5.13\times10^{-6}$ & $1.24\times10^{-2}$ & 5.95 \\
5  & 66.4  & $1.67\times10^{-2}$ & $5.12\times10^{-6}$ & $1.35\times10^{-1}$ & 8.11 \\
11 & 142.4 & $3.84\times10^{-3}$ & $3.15\times10^{-6}$ & $2.71\times10^{-2}$ & 7.04 \\
11 & 150.4 & $9.02\times10^{-3}$ & $3.15\times10^{-6}$ & $7.59\times10^{-2}$ & 8.42 \\
Healthy 0 & 200.0 & $5.91\times10^{-5}$ & $1.88\times10^{-7}$ & $1.96\times10^{-5}$ & 0.33 \\
Healthy 6 & 200.0 & $3.76\times10^{-4}$ & $1.56\times10^{-6}$ & $6.71\times10^{-4}$ & 1.78 \\
\bottomrule
\end{tabular}
\caption{Growth of the model response to its own state error.}
\label{tab:secant}
\end{table}

Spectrally, the induced response first appears in \(k=32,\ldots,64\), then
spreads through \(k=64,\ldots,128\).  This precedes the abrupt terminal event.
Thus the NaN is not born in one exceptional arithmetic operation; it is the
endpoint of a resolved spectral cascade.

The ratio \(q_n\) is a secant quantity, not an eigenvalue and not a proof that
the full Fr\'echet derivative has norm 8.  It is nevertheless the relevant
finite-amplitude response along the error direction actually selected by the
rollout.  The contrast with healthy cases matters more than its absolute
value.  In failing case 5 the induced response becomes
\(2.6\times10^4\) times the truth-state bias, while a healthy case remains at a
gain below one at final time.

The frame order provides a causal discriminator:
\begin{enumerate}
  \item The pointwise truth-state bias is small and nearly stationary.
  \item A small state discrepancy develops through ordinary accumulated error.
  \item The learned operator maps that discrepancy into a disproportionately
        large mid-band response.
  \item The response broadens spectrally and steepens the state.
  \item Only then do base-grid nonlinear products differ materially from their
        padded counterparts.
  \item Finally, the GL2 fixed-point residual ceases to contract and a
        nonfinite stage is produced.
\end{enumerate}
This ordering rules out ``one bad target row directly produces a NaN'' and
``the first cause is a single aliased product'' as descriptions of the observed
cases.  It does not claim that all possible future NaNs must follow the same
path; the conclusion is scoped to the remaining Tanaka population in the
audited checkpoints.

\section{Hamiltonian inconsistency is the root defect}

For a genuine DNO define
\begin{equation}
 B=\frac{G(\eta)\xi+\eta_x\xi_x}{1+\eta_x^2},
 \qquad V=\xi_x-B\eta_x.
\end{equation}
The Hadamard shape derivative satisfies, for every perturbation \(\zeta\),
\begin{equation}
 D_\eta G(\eta)[\zeta]\xi
 +G(\eta)(\zeta B)
 +\partial_x(\zeta V)=0.
 \label{eq:hadamard}
\end{equation}
This is an identity of operators applied to \(\xi\), not merely an energy
conservation statement along one trajectory.  Its importance can be seen by a
short variational calculation.  Let
\begin{equation}
 K_\theta(\eta,\xi)=\frac{1}{2}\langle\xi,\Gt(\eta)\xi\rangle,
\end{equation}
where \(\langle f,g\rangle=\int_0^{2\pi}fg\,\dd x\).  Assuming
self-adjointness of \(\Gt(\eta)\) and Eq.~\eqref{eq:hadamard}, periodic
integration by parts gives
\begin{align}
 \Deta K_\theta[\zeta]
 &=\frac12\inner{\xi}{\Deta\Gt(\eta)[\zeta]\xi} \\
 &=-\frac12\inner{\xi}{\Gt(\eta)(\zeta B_\theta)}
   -\frac12\inner{\xi}{\partial_x(\zeta V_\theta)} \\
 &=-\frac12\inner{\Gt(\eta)\xi}{\zeta B_\theta}
   +\frac12\inner{\xi_x}{\zeta V_\theta} \\
 &=\frac12\int_0^{2\pi}\zeta
       \left(V_\theta\xi_x-B_\theta\Gt(\eta)\xi\right)\dd x.
 \label{eq:kinetic-variation}
\end{align}
Therefore
\begin{equation}
 -\frac{\delta K_\theta}{\delta\eta}
 =\frac12\left(B_\theta\Gt(\eta)\xi-V_\theta\xi_x\right).
 \label{eq:kinetic-gradient}
\end{equation}
Using \(V_\theta=\xi_x-B_\theta\eta_x\) and
\(B_\theta=(\Gt\xi+\eta_x\xi_x)/(1+\eta_x^2)\), the right-hand side becomes
\begin{equation}
 -\frac12\xi_x^2
 +\frac12\frac{(\Gt(\eta)\xi+\eta_x\xi_x)^2}{1+\eta_x^2},
\end{equation}
which is exactly the nonlinear kinetic contribution used in
Eq.~\eqref{eq:zakharov-geometric}.  This proves the claimed equivalence.

Conversely, tying the learned Fourier multipliers makes \(\Gt\) self-adjoint
and guarantees \(\delta K_\theta/\delta\xi=\Gt\xi\), but it says nothing about
\(\delta K_\theta/\delta\eta\).  The surface derivative must separately obey
Eq.~\eqref{eq:hadamard}.  This is the precise sense in which self-adjointness is
necessary but not sufficient.

\subsection{The measured defect}

For diagnosis, the continuum derivative is evaluated by automatic
differentiation of the model energy and compared with the geometric expression
on the same discrete grid.  Both sides undergo the production
\(|k|\leq128\) projection.  Schematically, the reported relative defect is
\begin{equation}
 d_{\mathrm{shape}}
 =\frac{\norm{P_{\leq128}(
      \xi_{t,\mathrm{geom}}+\delta_\eta K_\theta
   )}_{H^1}}
 {\norm{P_{\leq128}\xi_{t,\mathrm{geom}}}_{H^1}+\varepsilon}.
 \label{eq:shape-defect}
\end{equation}
The same normalization, projection, and differentiation pipeline is used for
the order-six reference control.

The measured discrepancy, after the production projection to \(|k|\leq128\),
is summarized below.  The exact order-six DNO passes the same discrete test,
ruling out a diagnostic normalization error.

\begin{table}[ht]
\centering
\small
\begin{tabular}{lrrrr}
\toprule
Case/time & Value error & Defect at truth & Defect at rollout & Order-six defect \\
\midrule
Healthy 0 / 200.0 & $6.47\times10^{-5}$ & 0.0896 & 0.2966 & $2.0\times10^{-9}$ \\
5 / 33.6          & $2.34\times10^{-4}$ & 1.394  & 4.975  & $1.43\times10^{-5}$ \\
11 / 189.6        & $2.15\times10^{-4}$ & 0.502  & 1.117  & $3.43\times10^{-6}$ \\
22 / 45.6         & $5.54\times10^{-4}$ & 2.724  & 7.481  & $3.82\times10^{-6}$ \\
\bottomrule
\end{tabular}
\caption{Small pointwise error coexisting with a large shape-derivative defect.}
\label{tab:shape}
\end{table}

This resolves an apparent contradiction.  Gauss--Legendre methods are
symplectic for the Hamiltonian vector field supplied to them.  The surrogate
field is not, in general, the gradient of one Hamiltonian because its
\(\eta_t\) and \(\xi_t\) components use incompatible derivatives of \(\Gt\).
The time integrator therefore cannot protect the rollout from the learned
geometric inconsistency.

One can also phrase the problem through the error equation.  Write the truth
and surrogate vector fields as \(F\) and \(F_\theta\), and let
\(e=z_\theta-z\).  Before the error is large,
\begin{equation}
 e_t
 =\underbrace{F_\theta(z)-F(z)}_{\text{truth-manifold bias}}
 +\underbrace{DF_\theta(z)e}_{\text{propagation by learned tangent map}}
 +O(\norm{e}^2).
 \label{eq:error-linearization}
\end{equation}
Supervised one-step fitting directly attacks the first term.  The aligned NPZ
analysis shows that this term stays small while the second becomes dominant.
The Hadamard defect identifies a specific incorrect component of
\(DF_\theta\), rather than appealing vaguely to ``error compounding.''  The
spectral cascade is error compounding, but it is error compounding through a
geometrically wrong tangent operator.

\section{The missing spectral null form}

The exact first Craig--Sulem correction is
\begin{equation}
 G_1(\eta)\xi
 =-G_0\!\left(\eta G_0\xi\right)
  -\partial_x\!\left(\eta\partial_x\xi\right).
 \label{eq:g1}
\end{equation}
Its Fourier interaction kernel is
\begin{equation}
 \widehat{G_1(\eta)\xi}_k
 =\sum_{k'}
 \underbrace{\left(kk'-g_0(k)g_0(k')\right)}_{\mathcal K_1(k,k';h)}
 \widehat{\eta}_{k-k'}\widehat{\xi}_{k'},
 \qquad g_0(k)=|k|\tanh(h|k|).
 \label{eq:g1-kernel}
\end{equation}
The sign is worth checking.  The first term in Eq.~\eqref{eq:g1} contributes
\(-g_0(k)g_0(k')\), while
\(-\partial_x(\eta\partial_x\xi)\) contributes \(+kk'\).  In deep water,
\(g_0(k)=|k|\), so \(\mathcal K_1=0\) whenever \(k\) and \(k'\) have the same
sign.  This is not a learned small coefficient: it is exact cancellation
between two individually large derivative terms.  At finite depth the
cancellation is exponentially accurate when \(h|k|\) and \(h|k'|\) are large,
as they are near the observed modes.

The null form is physically important in the failing Tanaka regime.  A transfer-kernel probe
at a truth surface found that the learned model maps \(k_{\mathrm{in}}=117\)
to \(k_{\mathrm{out}}=115\) with amplitude 985.67, whereas the order-six
reference amplitude is \(1.47\times10^{-5}\), a ratio of approximately
\(6.7\times10^7\).  The flat diagonal response at this mode is about
\(5.9\times10^4\), so 985.67 is only about 1.7\% of the diagonal amplitude.
That explains why an aggregate pointwise error metric can appear excellent
while a repeatedly injected off-diagonal sideband matters dynamically.  The
learned model can represent the cancellation but does not enforce it.

Across all sampled Tanaka truth states, replacing \(G_0\) by exact \(G_0+G_1\)
reduces median relative \(G\xi\) error from 0.1496 to 0.00555 and p95 from
0.2089 to 0.01373.  More importantly, it makes the high-frequency null form
exact rather than dependent on cancellation among learned blocks.

The analytic-only rollout is a particularly discriminating experiment.  Exact
\(G_0+G_1\), with no cap and no learned residual, yields 0/32 NaNs and 0/32
divergences.  It nevertheless gives final \(\eta\) median 0.2071 and p95
0.4096.  Hence two statements are simultaneously true:
\begin{enumerate}
  \item the first-order null structure removes the catastrophic mechanism in
        this test; and
  \item neglected \(G_2,G_3,\ldots\) terms are essential for useful long-time
        phase and amplitude accuracy.
\end{enumerate}
This is why the proposed solution is not to abandon learning, but to prevent
the learned residual from overwriting the analytically known first-order term.

\section{Aliasing and the role of the integrator}

For truncated spectra, a physical-space product computes a circular
convolution.  If \(f\) and \(g\) contain modes near the retained cutoff, terms
with true wavenumber outside the grid wrap into lower modes.  The nonlinear
Zakharov right-hand side contains products and a rational expression in
\(\eta_x\), \(\xi_x\), and \(G\xi\), so this mechanism becomes severe once the
surrogate has populated its upper retained bands.

Xu and Guyenne explicitly identify aliasing in both the DNO expansion and the
equations of motion, and extend the spectra by a factor of two for nonlinear
terms in the equations of motion~\cite[Sec.~3.3]{xu2009}.  Our order-six
reference DNO already uses padding, but the original surrogate and truth
rollout paths evaluated the quadratic Zakharov products at base resolution.
The correction now implemented is:
\begin{enumerate}
  \item Fourier-interpolate \(\eta_x\), \(\xi_x\), and \(G\xi\) onto a grid
        with twice as many points;
  \item evaluate the nonlinear expression pointwise on that padded grid; and
  \item truncate the result back to the production spectrum before applying
        the usual retained-mode projection.
\end{enumerate}

Two-times Fourier padding is therefore a necessary discretization correction.
On healthy trajectories, base-grid and padded right-hand sides agree to about
\(10^{-11}\).  Near the last finite frames, the discrepancy becomes large:
\begin{table}[ht]
\centering
\small
\begin{tabular}{lrrrr}
\toprule
Case & Last finite time & Relative RHS difference & First $>10^{-3}$ & First $>0.1$ \\
\midrule
Healthy 0 & 200.0 & $2.71\times10^{-11}$ & never & never \\
Healthy 6 & 200.0 & $1.78\times10^{-11}$ & never & never \\
5 & 33.6 & 0.753 & 32.0 & 33.6 \\
11 & 189.6 & 0.819 & 187.2 & 188.8 \\
22 & 45.6 & 0.0829 & 45.6 & never \\
\bottomrule
\end{tabular}
\caption{Difference between base-grid and two-times-padded Zakharov right-hand
sides.}
\label{tab:rhs-alias}
\end{table}

The large terminal discrepancy alone does not establish that aliasing started
the cascade.  A causal replay does: restarting each failing state with padded
products delays the first nonfinite value by only 0.05, 0.08, and 0.05 time
units for cases 5, 11, and 22.  A healthy replay agrees to
\(3.5\times10^{-15}\).  Aliasing is therefore a selective terminal amplifier,
not the initiator in these cases.  The padding fix should be retained
regardless of C16's outcome because it restores the intended pseudospectral
discretization without perturbing resolved healthy trajectories.

Likewise, more Picard work or smaller steps cannot undo the preceding
operator-induced drift.  A residual check remains valuable as a reported
fail-safe, and adaptive halving may rescue milder future failures, but it is not
a substitute for restoring Hamiltonian compatibility.

There is no conflict with the well-known long-time advantages of symplectic
integration.  Those advantages concern the discretization of a Hamiltonian
vector field and usually imply bounded, structured energy error over long
times, rather than exact conservation at each step.  Here the learned vector
field fails the compatibility calculation in
Eq.~\eqref{eq:kinetic-variation}.  Moreover, the practical scheme uses a fixed
number of nonlinear iterations, so once the state enters a high-gain region it
may not even realize the exact implicit GL2 map.  The integrator is important,
but it is downstream of the operator geometry.

\section{C16: the resulting model and objective}

The proposed clean retraining uses
\begin{equation}
 \Gt(\eta)=G_0+G_1(\eta)+\R(\eta),
 \qquad \R(\eta)=O(\eta^2).
 \label{eq:c16}
\end{equation}
The analytic \(G_1\) path is uncut and evaluated with double-precision spectral
arithmetic.  The learned residual uses eight tied, self-adjoint blocks with 256
branches each.  A bias-free feature trunk followed by
\(f\mapsto f\tanh(f)\) ensures
\begin{equation}
 \R(0)\xi=0,
 \qquad D_\eta\R(0)[\zeta]\xi=0
\end{equation}
for every parameter value.  The residual sum is scaled by
\(1/\sqrt{8\cdot256}\), making the size of a fresh Adam update independent of
the chosen parallel rank without changing the function class.

\subsection{Residual parameterization and its invariants}

Suppressing depth and normalization factors, the residual has the schematic
form
\begin{equation}
 \R(\eta)\xi
 =\frac{1}{\sqrt{N_bN_r}}
  \sum_{b=1}^{N_b}\sum_{r=1}^{N_r}
  M_{b,r}(D,h)
  \left[\Phi_{b,r}(\eta)\,
        M_{b,r}(D,h)\xi\right],
 \quad N_b=8,\quad N_r=256.
 \label{eq:residual-block}
\end{equation}
The same real Fourier multiplier appears on both sides of the pointwise
multiplication.  For fixed \(\eta\), each term is self-adjoint:
\begin{align}
 \inner{\psi}{M[\Phi M\xi]}
 &=\inner{M\psi}{\Phi M\xi}
  =\inner{\Phi M\psi}{M\xi}
  =\inner{M[\Phi M\psi]}{\xi}.
\end{align}
The proof uses self-adjointness of the real multiplier and of multiplication by
the real function \(\Phi\).  Summation and positive scalar normalization
preserve the property.

For the order statement, let the bias-free eta trunk produce
\(f(\eta)=L\eta+O(\norm{\eta}^2)\) near the flat surface.  Since
\(\tanh f=f+O(f^3)\),
\begin{equation}
 f(\eta)\tanh f(\eta)
 =(L\eta)^2+O(\norm{\eta}^3).
\end{equation}
All subsequent eta projections are bias-free.  Consequently every
\(\Phi_{b,r}(0)=0\) and \(D_\eta\Phi_{b,r}(0)=0\), proving the two structural
zeros.  The bounded-growth lift was chosen instead of raw \(f^2\): it retains
quadratic behavior near zero but grows only linearly in \(|f|\) when derivative
features are large.  Learned mixing before the lift still permits general
quadratic cross-interactions; this is not a restriction to literal squares of
the physical input channels.

The architecture contains 860,928 trainable parameters.  It retains the full
eight-block capacity because the one-block and block-deletion experiments did
not support a capacity reduction.  There is no hard spectral cutoff, no output
cap, no PSD hinge, and no Tanaka-specific data pack in C16.  Thus a successful
result can be attributed to the null-structured backbone and direct shape
constraint rather than to spectral culling.

\subsection{Finite-secant Hadamard loss}

The randomized finite-secant residual is
\begin{align}
 S_\theta={}&
 \frac{\Gt(\eta+\epsilon\zeta)\xi-\Gt(\eta)\xi}{\epsilon}
 +\Gt(\eta)(\zeta B_\theta)
 +\partial_x(\zeta V_\theta),
\end{align}
and the training objective is
\begin{equation}
 \mathcal{L}
 =\mathcal{L}_{\mathrm{data}}
 +\lambda_{\mathrm{shape}}
 \mathbb{E}_{\zeta}
 \frac{\norm{W P_{\leq128}S_\theta}_2^2}
 {\norm{W P_{\leq128}\bigl(
     \Gt(\eta)(\zeta B_\theta)+\partial_x(\zeta V_\theta)
   \bigr)}_2^2+\delta}.
 \label{eq:loss}
\end{equation}
Here \(W\) is the supervised Sobolev weight.  Relative, zero-mean,
inverse-Sobolev probes cover low and mid/high modes without selecting specific
Tanaka cases.  The regularizer is evaluated in double precision on a global
microbatch of eight every 16 optimizer steps, with
\(\lambda_{\mathrm{shape}}=10^{-2}\) introduced through a warmup.

More explicitly, a random Fourier probe is projected to zero mean and scaled
by an inverse Sobolev symbol so its high modes do not dominate merely because
derivatives multiply by \(|k|\).  Its physical RMS is chosen relative to the
RMS of the current \(\eta\), subject to a small floor, and
\(\epsilon\in[10^{-3},3\times10^{-3}]\) is sampled as a relative secant size.
The two model evaluations use exactly the production input normalization,
output de-normalization, filtering, and output-mean subtraction.  Both the
secant residual and its normalizing forcing are measured after projection to
\(|k|\leq128\) in the same Sobolev metric used by supervised training.

The denominator in Eq.~\eqref{eq:loss} is essential.  An absolute residual
would overweight large-amplitude samples and permit the loss scale to change
with probe amplitude.  The relative form asks whether the identity error is
small compared with the physical terms that should cancel.  The floor
\(\delta\) only regularizes nearly trivial forcing.  A finite secant, rather
than an exact Jacobian norm, avoids differentiating through a full Jacobian and
keeps the cost compatible with sparse microbatch evaluation.

This loss is related to, but cleaner than, the earlier GL2 stage regularizer.
The stage loss asks whether one particular discretized short-time map behaves
well under selected perturbations; it mixes DNO error, time-discretization
details, filters, and stage construction.  The Hadamard loss asks directly for
the continuum operator identity that makes the Zakharov vector field
Hamiltonian.  It is therefore cheaper to interpret and easier to defend in a
mathematical presentation.

\subsection{Optimizer calibration}

The first implementation exposed a separate optimizer parameterization issue:
hundreds of thousands of zero-initialized projection weights received aligned
Adam sign updates and were summed without rank normalization.  Before the
normalization, a single real-state step at learning rate \(2\times10^{-5}\)
changed loss from 0.02463 to 0.3201.  After normalization, the same step reduced
it to 0.02174.  A 1\% v9 GPU smoke then completed with train loss 0.01361 and
validation loss 0.01311, with no surge or nonfinite metric.

The mechanism of the initial optimizer instability deserves emphasis because
it is independent of the NaN-rollout diagnosis.  At initialization, the
analytic baseline already has small loss and the per-block eta-to-branch
projection kernels are zero.  Adam normalizes each nonzero coordinate by its
own second moment, so the first update is approximately a sign step in each of
524,288 projection coordinates.  All coordinates are trained against the same
residual and their function-space effects add coherently across eight blocks
and 256 branches.  A parameter-space learning rate that was harmless in a
warm-started model therefore produced an enormous first function-space step in
the fresh model.

The factor \(1/\sqrt{N_bN_r}=1/\sqrt{2048}\) is a fan-in normalization for the
parallel residual paths.  It preserves exact \(G_0+G_1\) at initialization,
the order-two zero, self-adjointness, parameter count, and the set of functions
representable after rescaling downstream weights.  CPU calibration gave
\begin{center}
\begin{tabular}{lrr}
\toprule
Parameterization & Loss before step & Loss after AdamW step at $2\times10^{-5}$ \\
\midrule
Unnormalized order-two residual & 0.024634 & 0.320115 \\
Fan-in-normalized residual & 0.024634 & 0.021744 \\
\bottomrule
\end{tabular}
\end{center}
The data-gradient to Hadamard-gradient norm ratio was 0.0662, and the chosen
Hadamard weight 0.01 makes the auxiliary gradient influential but not dominant
at initialization.  Both the learning rate and Hadamard weight are warmed up
over 500 steps.

\subsection{What C16 guarantees, and what it does not}

For every value of its trained parameters, C16 guarantees:
\begin{itemize}
  \item exact \(G_0+G_1\) behavior through first order at \(\eta=0\);
  \item linearity and self-adjointness in \(\xi\) for the learned residual;
  \item no learned \(O(1)\) or \(O(\eta)\) term that can overwrite the analytic
        first-order null form; and
  \item a consistently two-times-dealiased Zakharov nonlinear right-hand side.
\end{itemize}
It does \emph{not} guarantee the full nonlinear Hadamard identity for finite
\(\eta\); that is encouraged statistically by Eq.~\eqref{eq:loss}.  Nor does it
prove bounded rollouts, exact energy conservation, or uniform operator error
outside the sampled distribution.  Calling it an ``explicit bounded residual
architecture'' would therefore be inaccurate.  Its principle is geometric
compatibility and asymptotic null structure, not a hard global norm bound.

\section{Current experiment and acceptance criteria}

At the time of writing, a budget-capped C16 retraining is running on two H100
GPUs in the personal Modal workspace.  It uses clean v9 data, batch size 1024,
learning rate \(2\times10^{-5}\), 500-step learning-rate and Hadamard warmups,
and 16 epochs.  The abbreviated run is intended to obtain the strongest
checkpoint possible within approximately 30 dollars of workspace credit; it is
not assumed that epoch 16 is fully converged.

The first two complete epochs are healthy:
\begin{center}
\begin{tabular}{lrrrr}
\toprule
Epoch & Train loss & Validation loss & Hadamard loss & Defect RMS \\
\midrule
1 & 0.009919 & 0.008740 & 0.04381 & 0.1700 \\
2 & 0.008472 & 0.008376 & 0.03104 & 0.1482 \\
\bottomrule
\end{tabular}
\end{center}
The regularizer was active on 373 of 5964 batches in each epoch, matching the
intended frequency \(1/16\).  Its mean effective weight reached
\(9.57\times10^{-3}\) in epoch 1 and the full \(10^{-2}\) in epoch 2.  These
are encouraging training diagnostics, not rollout evidence.  In particular,
the random-probe defect RMS cannot yet be compared directly with the
case-specific defect table without running the same held-state diagnostic.

An earlier Modal attempt stopped one third of the way through epoch 1 because
the local monitoring client was explicitly interrupted; Modal logs show a
cancellation signal rather than a Python, JAX, or model exception.  No
checkpoint was produced by that attempt.  The current 16-epoch run is a fresh
relaunch and should be judged independently.

The structural hypothesis should be accepted only if rollout evaluation shows
all of the following:
\begin{enumerate}
  \item zero Tanaka NaNs and divergences among 32 truth-valid cases;
  \item zero scored Benjamin--Feir NaNs and divergences among 31 cases;
  \item no material median or p95 regression relative to the strongest prior
        checkpoints;
  \item at least a two-order-of-magnitude reduction in the failing-state shape
        defect;
  \item no pre-cascade secant-gain growth toward 7--8; and
  \item healthy GL2 fixed-point residuals without relying on output caps,
        hard spectral cutoffs, or case-specific damping.
\end{enumerate}

The main scientific conclusion does not depend on the outcome of one training
run: long-rollout stability requires approximating the geometry of the DNO in
addition to its values.  C16 is the first experiment in this sequence that
directly enforces the missing first-order null form and shape derivative while
retaining the capacity needed for higher-order accuracy.

\subsection{Planned post-training decision sequence}

To avoid another ambiguous experiment, evaluation should proceed in the
following order:
\begin{enumerate}
  \item Verify checkpoint/config integrity and run the held-state Hadamard
        defect table on the final checkpoint.  This is cheap and directly tests
        whether the intended loss changed the diagnosed quantity.
  \item Evaluate the final checkpoint on Tanaka \(g_0\) and Benjamin--Feir
        \(g_1\) with cached truth, the double-precision harness, batched
        surrogate execution, and no run-time limiter or damping.
  \item If final is negative, evaluate best-validation only as a checkpoint
        selection control; do not reinterpret it as a separate architecture.
  \item For every remaining failure, regenerate the frame-aligned secant,
        spectral-band, RHS-alias, and GL2-residual traces.  The question is not
        merely whether the failure count moved, but whether the causal signature
        changed.
  \item Only after the two decisive regimes pass should the full multi-regime
        suite be run to rule out regressions in shallow water, random seas,
        Stokes waves, and modal instability.
\end{enumerate}

\section{Alternative explanations, objections, and falsifiers}

This section states the questions most likely to arise in a rigorous oral
discussion and the limits of the current answers.

\subsection{Could the data still contain subtle corruption?}

Yes, in the weak logical sense that no finite audit proves every stored number
correct.  What has been ruled out is the form of corruption needed to explain
the observed failure mechanism.  Random and source-stratified scans found no
nonfinite base rows; sign and spectral heuristics did not reveal a population
aligned with the failing cases; and, most importantly, the learned DNO has
small value error when evaluated on the exact failing truth states.  A corrupted
label population severe enough to seed the reported sideband directly should
normally appear as increased error on those states.  Instead, the large error
appears when the input is moved in the rollout's transverse direction.

A useful remaining audit would compare independently regenerated order-six
labels against stored v9 labels on a stratified sample, including source,
depth, steepness, and high-band-energy tails.  Such an audit would strengthen
dataset provenance but would not replace the already measured shape identity.

\subsection{Does small value error plus large derivative error necessarily
cause a NaN?}

No.  A large local derivative defect is a mechanism and risk factor, not by
itself a theorem of finite-time blow-up.  The causal case is empirical: the
defect is present at the failing states, the realized secant direction grows,
its spectrum matches the subsequent cascade, and analytic restoration of the
first-order null form removes all NaNs in the 32-case test.  A formal theorem
would require stability bounds for the coupled semi-discrete system and a
uniform neighborhood estimate for \(D\Gt\), neither of which is currently
available.

\subsection{Could GL2 still be the primary issue?}

The fixed number of Picard iterations is certainly where nonfiniteness is
finally emitted, and a more robust nonlinear solver could extend some
trajectories.  It is unlikely to be the primary issue in the audited cases
because the model-induced mid-band response is already orders of magnitude
above its truth-state bias before the stage residual fails; substep refinement
and padded-RHS replay delay failure only slightly; and healthy cases use the
same integrator without trouble.  A definitive follow-up could solve selected
near-terminal stages with Newton--Krylov to high tolerance.  If the converged
implicit map returned to a healthy spectral state, the present conclusion
would need revision.  If it advanced the already contaminated state, it would
confirm the current ordering.

\subsection{Why not learn the whole Hamiltonian and differentiate it?}

That is a principled alternative.  A scalar Hamiltonian network with
\(\eta_t=\delta_\xi H_\theta\) and
\(\xi_t=-\delta_\eta H_\theta\) would be Hamiltonian by construction.  The
tradeoff is computational: rollout requires derivatives of a nonlocal learned
energy, training DNO values can involve mixed derivatives, and encoding the
known linear dispersion and Craig--Sulem structure is less direct.  C16 keeps
the efficient DNO interface while enforcing the identity that makes that
interface Hamiltonian.  If C16 fails despite a substantially reduced measured
shape defect, a direct energy architecture becomes the most principled next
architectural family.

\subsection{Why not impose a global Lipschitz or spectral norm bound?}

A uniform bound could control amplification but would not distinguish physical
from unphysical high-frequency response.  The DNO is an order-one nonlocal
operator; its norm naturally scales with wavenumber, and nonlinear wave
dynamics legitimately transfer energy across modes.  A naive global Lipschitz
constant or hard cutoff risks destroying exactly the dynamics we are trying to
approximate.  The Hadamard identity is selective: it constrains how the large
terms must cancel, rather than requiring every term to be small.

\subsection{What result would falsify the C16 mechanism?}

The strongest falsifier would be: the final C16 checkpoint reduces the held
shape defect by two orders of magnitude and preserves the same-sign \(G_1\)
null form, yet the same cases still develop the same pre-cascade secant gain
and spectral signature.  That would show that the diagnosed identity is not
the controlling mechanism.  Other negative outcomes are less decisive:
\begin{itemize}
  \item If the shape defect does not improve, the loss estimator, weight, or
        16-epoch budget failed to enforce the hypothesis; the hypothesis itself
        was not tested strongly.
  \item If the defect improves but one-step accuracy is poor, the structural
        constraint may be correct but the residual architecture or optimization
        is inadequate.
  \item If NaNs disappear but p95 error becomes unacceptable, C16 is a stability
        result but not a useful surrogate.
\end{itemize}

\appendix
\section{Chronological experiment map}

This appendix is deliberately redundant with \texttt{EXPERIMENTS.md}.  Its
purpose is to make the logic reconstructible without memorizing labels.

\begin{longtable}{@{}p{0.10\linewidth}p{0.22\linewidth}p{0.29\linewidth}p{0.30\linewidth}@{}}
\toprule
Handle & Intervention & Outcome & What it taught us \\
\midrule
\endfirsthead
\toprule
Handle & Intervention & Outcome & What it taught us \\
\midrule
\endhead
v3--v10 & Full-feature CS-DNO baselines; progressively broader clean data and a
double-precision rollout harness. & Much better one-step and median rollout
accuracy than generic FNOs, but persistent rare Tanaka NaNs (v10: 4/32). & The
Craig--Sulem inductive bias is valuable, but pointwise accuracy and broad data
coverage do not ensure rollout stability. \\

C1 & Light GL2 stage-tangent matching fine-tune. & Tanaka 4/32 to 3/32 NaNs;
Benjamin--Feir 2/31 to 0/31. & Off-manifold sensitivity is a real lever, though
the indirect proxy is incomplete. \\

C2, first launch & Stage match plus gain, but resumed a tied checkpoint into an
untied architecture. & 32/32 NaNs in both decisive regimes. & Invalid
experiment; motivated strict resume architecture and parameter-count guards. \\

C2, corrected & Same intended loss with tied architecture preserved; final
checkpoint evaluated after the auxiliary loss acted. & Tanaka 3/32 NaNs,
median 0.00907; BF 0/31, median 0.00735. & Final versus best checkpoint matters
for auxiliary-loss fine-tunes; still short of zero NaNs. \\

GL2 guard & Picard residual detection, attempted halving, then per-sample
masking. & Detected but did not reduce NaN counts; early fleet-wide logic could
couple one bad sample to the whole batch. & GL2 noncontractivity is downstream;
batch-local safety logic must remain per sample. \\

Block-7 ablation & Zero an apparently amplifying learned block. & Tanaka NaNs
rose from 4/32 to 28/32; median error rose 62-fold. & Trained blocks form a
fragile compensating sum; post hoc deletion is not an architectural study. \\

v11-D & 80-fold smaller model plus cutoff, stage regularizer, and PSD hinge. &
Same Tanaka NaN count as baseline and 12--20-fold median regressions. & A
kitchen-sink run is hard to interpret and reduced capacity was not a cure. \\

v11-E & Full capacity with hard learned-block \(k\)-cut isolated. & Did not
establish a useful zero-NaN/accuracy tradeoff. & Brute-force band removal is
not the preferred mechanism. \\

C6/C7 & Model-harvested pre-cascade hard negatives. & Harmful or non-general
fine-tunes. & Contaminated off-manifold packs can feed the model's own failure
back into training. \\

C8/C10 & Increasingly matched stage-gain regularization, with C10 using the
production filter and \(k=64,\ldots,128\) target band. & C10 reached the best
pre-C16 count, 2/32 Tanaka and 0/31 BF NaNs, but regressed some medians/tails. &
The diagnosed spectral band matters; indirect regularization still does not
restore the defining operator identity. \\

C9 & Stronger soft full-output high-band envelope penalty. & Penalty diagnostic
barely moved; Tanaka remained 3/32 NaNs and accuracy regressed. & Simply
increasing this penalty weight is not an effective axis. \\

C10 + runtime limiter & Cap learned \(G\xi\) high-band output during rollout. &
Tanaka stayed 2/32 NaNs; BF p95 worsened. & The operational guard does not reach
zero and can damage tails. \\

C11/C13 & Clean truth-generated hard-negative packs near selected remaining
failures. & Training was feasible; C13 evaluation was superseded before a final
claim. & Useful as a local learnability probe, not a universal or satisfying
paper-level mechanism. \\

C12 & One-step pushforward/filtering fallback prepared from C2. & Not promoted
after the stronger structural diagnosis. & Rollout-state training remains an
option, but it mixes integrator and operator effects. \\

C14 & One-block, checkpoint-carrying full-output cap. & Best finite checkpoint:
3/32 Tanaka NaNs, 6/32 divergences, median 0.162, p95 1.253; training later
became nonfinite. & The cap has a mid-band/floor loophole and severe accuracy
cost; one-block capacity was also a confounder. \\

C15 & Eight-block capacity control for the C14 cap. & Interrupted before epoch
1 and before any checkpoint. & Inconclusive as a capacity ablation, but not
worth relaunching because the cap mechanism itself was analytically rejected. \\

Root-cause audit & Frame alignment, realized secants, automatic-differentiation
shape check, transfer kernel, analytic \(G_0+G_1\), and padded-RHS replay. &
Small truth-state value error coexists with \(O(1)\) shape defect; exact
\(G_0+G_1\) is stable; RHS padding delays rather than cures failure. & The DNO
shape derivative/Hamiltonian compatibility is the primary missing constraint. \\

C16 & Fresh exact \(G_0+G_1\), full self-adjoint \(O(\eta^2)\) residual,
randomized Hadamard loss, and padded RHS. & Structural tests and 1\% smoke pass;
16-epoch budget-capped training is in progress. & First clean direct test of the
diagnosis. \\
\bottomrule
\end{longtable}

\section{Implementation and correctness checks}

The implementation was not accepted on training loss alone.  The following
independent checks passed before the full launch:
\begin{itemize}
  \item five architecture tests: legacy order-one identity, zero order-two
        value and first variation, quadratic near-zero scaling,
        self-adjointness, and isolation of double precision to analytic \(G_1\);
  \item seven Hadamard-loss tests: projected Sobolev energy, probe scaling and
        bandwidth, exact flat-surface identity, detection of missing \(G_1\),
        production normalization/mean behavior, JIT and mixed-precision
        gradients, and invalid-configuration handling;
  \item two dedicated padded-RHS tests: equality on bandlimited products and
        removal of high-frequency quadratic aliasing;
  \item fourteen existing stage-regularizer tests after routing all nonlinear
        right-hand sides through the same padded helper; and
  \item a 59-step two-GPU smoke with finite checkpoints and validation.
\end{itemize}

\section*{Reproducibility record}

The detailed chronological record is \texttt{EXPERIMENTS.md}.  Frame-level
statistics and paths to the aligned NPZ archives are recorded in
\texttt{notes/nan\_root\_cause\_investigation\_20260709.md}.  The C16 launcher
is \texttt{scripts/launch\_c16\_modal\_h100x2.sh}; model, Hadamard-loss, and
dealiased-RHS invariant tests are stored alongside their implementations.

\begin{thebibliography}{9}
\bibitem{xu2009}
L. Xu and P. Guyenne,
``Numerical simulation of three-dimensional nonlinear water waves,''
\emph{Journal of Computational Physics}, vol.~228, pp.~8446--8466, 2009,
doi:10.1016/j.jcp.2009.08.015.
\end{thebibliography}

\end{document}
